% =================================================================
% Rapport — Projet de Recherche Opérationnelle
% Algorithmes de flot (Ford-Fulkerson, MCF, détection de cycle négatif)
%
% Compilation : pdflatex rapport.tex  (deux passes pour la TOC)
% Compatible Overleaf.
% =================================================================

\documentclass[11pt,a4paper]{article}

\usepackage[utf8]{inputenc}
\usepackage[T1]{fontenc}
\usepackage[french]{babel}
\usepackage{lmodern}

\usepackage[a4paper, margin=2.5cm]{geometry}
\usepackage{microtype}
\usepackage{parskip}
\setlength{\parindent}{0pt}

\usepackage{amsmath, amssymb}

\usepackage{xcolor}
\definecolor{accent}{HTML}{2A6F97}
\definecolor{softgray}{HTML}{F4F4F4}

\usepackage{algorithm}
\usepackage{algpseudocode}
\algrenewcommand\algorithmicwhile{\textbf{tant que}}
\algrenewcommand\algorithmicdo{\textbf{faire}}
\algrenewcommand\algorithmicend{\textbf{fin}}
\algrenewcommand\algorithmicif{\textbf{si}}
\algrenewcommand\algorithmicthen{\textbf{alors}}
\algrenewcommand\algorithmicelse{\textbf{sinon}}
\algrenewcommand\algorithmicfor{\textbf{pour}}
\algrenewcommand\algorithmicreturn{\textbf{renvoyer}}
\floatname{algorithm}{Algorithme}

\usepackage{listings}
\lstset{
    basicstyle=\ttfamily\small,
    keywordstyle=\color{accent}\bfseries,
    commentstyle=\color{gray}\itshape,
    showstringspaces=false,
    breaklines=true,
    frame=single,
    framesep=4pt,
    rulecolor=\color{gray!40},
    backgroundcolor=\color{softgray},
    language=Python
}

\usepackage{tikz}
\usetikzlibrary{arrows.meta, positioning}

\usepackage{booktabs}
\usepackage[colorlinks=true, linkcolor=accent, urlcolor=accent]{hyperref}

\newcommand{\Gres}{G_{r}}

\title{
    \vspace{-1.5cm}
    \textbf{Projet de Recherche Opérationnelle} \\[0.4em]
    \large Algorithmes de flot dans les réseaux
}
\author{Étudiant M1 --- UNICA}
\date{Année universitaire 2025--2026}

\begin{document}

\maketitle
\vspace{-1em}

\tableofcontents
\vspace{1em}
\hrule
\vspace{1em}

% =================================================================
\section{Introduction}

Ce projet implémente en Python pur (sans bibliothèque externe spécialisée)
trois familles d'algorithmes de flot :
\begin{itemize}
    \item flot maximum par \textbf{Ford--Fulkerson} (variante
        \emph{Edmonds--Karp}, BFS) ;
    \item flot de coût minimum (MCF) par augmentations successives, en deux
        versions : \textbf{Bellman--Ford} et \textbf{Dijkstra + Johnson} ;
    \item détection de \textbf{cycles négatifs} (Bellman--Ford, fallback DFS).
\end{itemize}
Seules \texttt{collections} et \texttt{heapq} sont utilisées.

% =================================================================
\section{Structure du graphe résiduel}
\label{sec:residuel}

\subsection{Choix de la représentation}

Deux options classiques :
\begin{itemize}
    \item \emph{matrice d'adjacence} : accès $O(1)$ mais $\Theta(V^2)$ en
        mémoire, parcours des voisins en $O(V)$ ;
    \item \emph{dict de dict} : \texttt{graph[u][v] =
        \{capacity, flow, cost\}}, accès $O(1)$ amorti, mémoire $O(|E|)$,
        parcours en $O(\deg(u))$.
\end{itemize}
On a retenu le \textbf{dict de dict}, plus adapté aux graphes creux qu'on
manipule en MCF.

\subsection{Astuce du flot symétrique}

Pour chaque arc réel $(u, v)$ ajouté, on crée \emph{aussi} une entrée
$\texttt{graph}[v][u]$ représentant l'arc inverse résiduel (capacité~$0$,
coût~$-w$) :

\begin{center}
\begin{tikzpicture}[->, >={Stealth[length=2.5mm]}, thick, node distance=4cm]
    \tikzstyle{node}=[circle, draw=accent, fill=accent!10, minimum size=1cm, font=\bfseries]
    \node[node] (u) {u};
    \node[node, right=of u] (v) {v};
    \draw[bend left=20, accent] (u) to node[above, font=\small, midway]
        {cap=5,~flow=3,~cost=2} (v);
    \draw[bend left=20, gray!70, dashed] (v) to node[below, font=\small, midway]
        {cap=0,~flow=$-3$,~cost=$-2$} (u);
\end{tikzpicture}
\end{center}

Le flot est stocké de manière \emph{symétrique} (positif sur l'arc réel,
négatif sur l'inverse) pour que la formule unique
\[
    \text{residual\_capacity}(u, v) = \text{capacity}(u, v) - \text{flow}(u, v)
\]
fonctionne dans les deux directions. Ainsi, augmenter de $\delta$ sur
$(u, v)$ se réduit à :
\begin{lstlisting}
graph[u][v]['flow'] += delta
graph[v][u]['flow'] -= delta
\end{lstlisting}
Cette uniformité évite de manipuler deux structures séparées et rend le
code des algorithmes beaucoup plus propre.

\subsection{Limitation}

Les arcs antiparallèles réels ($(u,v)$ et $(v,u)$ tous deux réels) ne sont
pas supportés par cette structure --- on lève une \texttt{ValueError}
plutôt que d'accepter silencieusement un état incohérent. La solution
standard consisterait à dédoubler l'un des deux arcs avec un sommet
intermédiaire.

% =================================================================
\section{Flot maximum : Edmonds--Karp}
\label{sec:ff}

\subsection{Principe}

Tant qu'il existe un chemin augmentant $s \to t$ dans le graphe résiduel,
on pousse le bottleneck le long de ce chemin et on met à jour le résiduel.
On choisit le chemin par \textbf{BFS} (plus court en nombre d'arcs), ce
qui garantit la borne polynomiale d'Edmonds--Karp.

\begin{algorithm}[H]
\caption{Edmonds--Karp}
\begin{algorithmic}[1]
\Function{FordFulkerson}{$G$, $s$, $t$}
    \State $\text{maxflow} \gets 0$
    \While{vrai}
        \State $P \gets \Call{BFS}{\Gres, s, t}$
        \If{$P = \emptyset$} \textbf{sortir} \EndIf
        \State $\delta \gets \min \{c_r(u,v) : (u,v) \in P\}$
        \State Augmenter $P$ de $\delta$
        \State $\text{maxflow} \mathrel{+}= \delta$
    \EndWhile
    \State $S \gets$ sommets atteignables depuis $s$ dans $\Gres$
    \State $T \gets V \setminus S$
    \State \Return $(\text{maxflow},\ S,\ T,\ \text{coupe} = \{(u,v) \text{ réel} : u \in S, v \in T\})$
\EndFunction
\end{algorithmic}
\end{algorithm}

\subsection{Complexité}

\[
\boxed{\;O(V \cdot E^2)\;}
\]
(BFS en $O(V+E)$, $O(VE)$ augmentations par le résultat d'Edmonds--Karp.)

\subsection{Coupe minimum}

Après convergence, un BFS depuis $s$ sur le résiduel donne l'ensemble $S$
des sommets atteignables. Les arcs réels $(u, v)$ avec $u \in S$ et
$v \notin S$ forment la coupe minimum (ils sont saturés). Par le théorème
max-flow / min-cut, la capacité de cette coupe est égale à la valeur du
flot maximum --- propriété utilisée comme sanity check dans les tests.

% =================================================================
\section{MCF par Bellman--Ford}
\label{sec:mcf-bf}

\subsection{Principe}

Analogue à Ford--Fulkerson, mais on choisit à chaque itération le chemin
de \emph{coût minimum} dans le résiduel (au lieu du plus court en
nombre d'arcs). On utilise \textbf{Bellman--Ford} parce que le résiduel
contient des arcs inverses de coût $-w$, donc Dijkstra ne s'applique pas
directement.

\subsection{Pré-requis : pas de cycle négatif}

Si le graphe contient un cycle de coût strictement négatif et de capacité
positive, le problème est \emph{non borné} (on peut faire tourner du flot
autour du cycle indéfiniment). On lance donc une détection préalable et
on avorte avec un message clair en cas de cycle.

\subsection{Pseudo-code}

\begin{algorithm}[H]
\caption{MCF par augmentations successives (Bellman--Ford)}
\begin{algorithmic}[1]
\Function{MCF-BF}{$G$, $s$, $t$, $F_{\max}$}
    \If{$G$ contient un cycle négatif} \textbf{abort} \EndIf
    \State $\text{cost}, \text{flow} \gets 0, 0$
    \While{$\text{flow} < F_{\max}$}
        \State $(d, \text{parent}) \gets \Call{Bellman-Ford}{\Gres, s}$
        \If{$d[t] = +\infty$} \textbf{sortir} \EndIf
        \State $P \gets$ chemin reconstruit
        \State $\delta \gets \min\{c_r(u,v) : (u,v) \in P\}$
        \State $\text{cost} \mathrel{+}= d[t] \cdot \delta$, \quad $\text{flow} \mathrel{+}= \delta$
        \State Augmenter $P$ de $\delta$
    \EndWhile
    \State \Return $(\text{cost}, \text{flow})$
\EndFunction
\end{algorithmic}
\end{algorithm}

\subsection{Complexité}

\[
\boxed{\;O(F \cdot V \cdot E)\;}
\]
où $F$ est le volume de flot poussé. Lourd dès que $F$ est grand, d'où
la variante Dijkstra+Johnson.

% =================================================================
\section{MCF par Dijkstra + Johnson}
\label{sec:mcf-johnson}

\subsection{Idée}

On veut Dijkstra à chaque itération (rapide grâce à \texttt{heapq}) mais
le résiduel contient des coûts négatifs. La solution de Johnson est de
\emph{reweighter} les coûts à l'aide de \textbf{potentiels} $h : V \to \mathbb{R}$ :
\[
    w'(u, v) := w(u, v) + h[u] - h[v].
\]

Deux propriétés essentielles de cette transformation :
\begin{enumerate}
    \item \textbf{Préservation des plus courts chemins} : le long d'un
        chemin $s \to t$, les termes $h[u] - h[v]$ se télescopent et la
        somme totale ne change que d'une constante ($h[s] - h[t]$). Le
        chemin optimal pour $w'$ est donc aussi optimal pour $w$.
    \item \textbf{Positivité des coûts reweightés} : si $h$ est choisi
        comme les distances depuis $s$ dans le résiduel, alors
        $w'(u, v) \ge 0$ partout (conséquence de l'inégalité triangulaire
        des plus courtes distances).
\end{enumerate}

\subsection{Initialisation et mise à jour}

\begin{itemize}
    \item Initialement, $h$ est calculé par \textbf{un seul} Bellman--Ford
        sur le graphe original.
    \item Après chaque Dijkstra, on met à jour $h[v] \mathrel{+}= d[v]$, ce
        qui maintient la positivité de $w'$ pour l'itération suivante.
\end{itemize}

\subsection{Récupération du coût réel}

Comme Dijkstra travaille sur $w'$, on obtient le coût réel par la formule
de télescopage :
\[
    \boxed{\;\text{cost réel}(s \to t) = d_{w'}[t] + h[t] - h[s]\;}
\]

\subsection{Pseudo-code}

\begin{algorithm}[H]
\caption{MCF par Dijkstra + reweighting de Johnson}
\begin{algorithmic}[1]
\Function{MCF-Johnson}{$G$, $s$, $t$, $F_{\max}$}
    \If{$G$ contient un cycle négatif} \textbf{abort} \EndIf
    \State $h \gets \Call{Bellman-Ford}{G, s}$ \Comment{potentiels initiaux}
    \State $\text{cost}, \text{flow} \gets 0, 0$
    \While{$\text{flow} < F_{\max}$}
        \State $(d, \text{parent}) \gets \Call{Dijkstra}{\Gres, s, w'}$
        \If{$d[t] = +\infty$} \textbf{sortir} \EndIf
        \State $\text{cost\_unit} \gets d[t] + h[t] - h[s]$
        \State $P \gets$ chemin reconstruit
        \State $\delta \gets \min\{c_r(u,v) : (u,v) \in P\}$
        \State $\text{cost} \mathrel{+}= \text{cost\_unit} \cdot \delta$, \quad $\text{flow} \mathrel{+}= \delta$
        \State Augmenter $P$ de $\delta$
        \State $h[v] \mathrel{+}= d[v]$ pour tout $v$ atteignable
    \EndWhile
    \State \Return $(\text{cost}, \text{flow})$
\EndFunction
\end{algorithmic}
\end{algorithm}

\subsection{Complexité}

\[
\boxed{\;O\bigl(V \cdot E + F \cdot (V + E) \log V\bigr)\;}
\]
Strictement meilleur que la version Bellman--Ford pure.

% =================================================================
\section{Détection de cycles négatifs}
\label{sec:neg-cycle}

\subsection{Méthode principale : Bellman--Ford}

Après $|V|-1$ itérations de relaxation, Bellman--Ford a stabilisé les
distances \emph{si et seulement si} il n'y a pas de cycle négatif
accessible. Une $|V|$-ième itération qui améliore encore une distance
signale donc l'existence d'un cycle négatif.

Pour détecter les cycles non accessibles depuis une source donnée, on
initialise $\text{dist}[v] = 0$ pour tout $v$ (équivalent à une source
virtuelle reliée à tous les sommets).

\subsection{Extraction du cycle}

Quand un sommet $v$ se fait relaxer à la $|V|$-ième itération, on remonte
$|V|$ fois via le tableau \texttt{parent} (garantissant qu'on aboutit
dans le cycle), puis on suit \texttt{parent} jusqu'à retomber sur le
sommet de départ pour reconstituer le cycle.

\subsection{Fallback DFS}

On a aussi une variante DFS qui trouve un cycle quelconque (via détection
d'arête arrière) puis vérifie son poids. Plus limitée que Bellman--Ford
mais utile comme second avis sur les petites instances.

% =================================================================
\section{Comparaison des deux approches MCF}
\label{sec:comparaison}

\begin{center}
\renewcommand{\arraystretch}{1.35}
\begin{tabular}{>{\bfseries}p{4.0cm} p{4.7cm} p{5.0cm}}
\toprule
\textbf{Critère} & \textbf{MCF Bellman--Ford} & \textbf{MCF Dijkstra + Johnson} \\
\midrule
Complexité / itération    & $O(V \cdot E)$               & $O((V+E)\log V)$ \\
Complexité totale         & $O(F \cdot V \cdot E)$       & $O(VE + F(V+E)\log V)$ \\
Coûts négatifs            & Directement                  & Après reweighting initial \\
Difficulté du code        & Faible                       & Moyenne \\
\bottomrule
\end{tabular}
\end{center}

\textbf{En pratique :} sur nos jeux de tests (4 à 10 sommets), les deux
algorithmes finissent en quelques millisecondes ; la différence devient
visible sur les instances plus grosses ou à grandes capacités. La
\textbf{cross-validation} entre les deux (test
\texttt{test\_mcf\_dijkstra\_johnson\_matches\_bf}) est la meilleure
\emph{sanity check} de l'implémentation : on vérifie qu'ils renvoient
exactement le même coût et le même flot.

% =================================================================
\section{Tests et résultats}

Le fichier \texttt{tests/test\_all.py} contient 11 tests couvrant :
\begin{itemize}
    \item flot max + conservation du flot + théorème max-flow / min-cut ;
    \item MCF Bellman--Ford avec arc de coût négatif ;
    \item MCF Dijkstra+Johnson, comparé à MCF-BF sur les mêmes graphes ;
    \item détection de cycle négatif (vrai positif et vrai négatif) ;
    \item avortement propre de MCF face à un cycle négatif ;
    \item stress test sur 10 sommets.
\end{itemize}
Tous les tests passent (\texttt{11/11}). L'exécution de
\texttt{main.py} produit les sorties demandées dans le sujet.

% =================================================================
\section{Conclusion}

Trois enseignements principaux :
\begin{enumerate}
    \item Le choix de la structure de données est central : l'astuce du
        \textbf{flot symétrique} (flot négatif sur l'arc inverse) unifie
        le traitement des deux directions et simplifie tout le code en aval.
    \item La \textbf{renormalisation de Johnson} permet de remplacer
        Bellman--Ford par Dijkstra dans le MCF, avec un gain de
        complexité significatif sur les instances de taille réaliste.
    \item La \textbf{cross-validation} de deux implémentations
        indépendantes (Bellman--Ford vs Dijkstra+Johnson) reste l'outil
        de validation le plus convaincant.
\end{enumerate}

\end{document}
